\ulohahead{společná matematika \textnormal{\footnotesize[úroveň 3]}}{Posloupnosti, limity a řady}
\begin{zadani}
\begin{enumerate}[label=(\alph*)]
  \item \bod{1} Definujte limitu posloupnosti a konvergenci. Uveďte aritmetiku limit a větu o dvou policajtech.
  \item \bod{2} Spočtěte $\displaystyle\lim_{n\to\infty}\frac{2n^2+1}{n^2+n}$ a $\displaystyle\lim_{n\to\infty}\sqrt[n]{n}$.
  \item \bod{3} Definujte součet řady a integrálním kritériem rozhodněte o konvergenci. Proč harmonická řada diverguje a geometrická (pro $|q|<1$) konverguje?
\end{enumerate}
\end{zadani}

\begin{reseni}
\textbf{(a)} $\lim_{n\to\infty}a_n=L$, pokud $\forall\varepsilon>0\,\exists n_0\,\forall n\ge n_0:|a_n-L|<\varepsilon$ (posloupnost je pak \emph{konvergentní}). \emph{Aritmetika limit}: limita součtu/součinu/podílu je součet/součin/podíl limit (u podílu při nenulové limitě jmenovatele). \emph{Věta o dvou policajtech}: je-li $a_n\le b_n\le c_n$ a $\lim a_n=\lim c_n=L$, pak $\lim b_n=L$.

\textbf{(b)} $\lim\frac{2n^2+1}{n^2+n}=\lim\frac{2+1/n^2}{1+1/n}=\frac{2}{1}=2$ (vydělením $n^2$). $\lim\sqrt[n]{n}=1$ (např.\ $\sqrt[n]{n}=e^{\frac{\ln n}{n}}$ a $\frac{\ln n}{n}\to0$).

\textbf{(c)} \emph{Součet řady} $\sum a_n$ je limita částečných součtu $s_N=\sum_{n=1}^N a_n$ (pokud existuje). \emph{Integrální kritérium} (pro kladnou klesající $f$): $\sum_{n\ge1} f(n)$ konverguje právě tehdy, když konverguje $\int_1^\infty f(x)\,dx$. \emph{Harmonická} $\sum\frac1n$ diverguje ($\int_1^\infty\frac{dx}{x}=\infty$; částečné součty rostou jako $\ln N$); \emph{geometrická} $\sum_{n=0}^\infty q^n$ pro $|q|<1$ konverguje k $\frac{1}{1-q}$ (částečný součet $\sum_{n=0}^N q^n=\frac{1-q^{N+1}}{1-q}$; součet od $n=1$ je $\frac{q}{1-q}$). \emph{(Nad rámec požadavku existují i srovnávací, podílové a odmocninové kritérium.)}
\end{reseni}

\begin{jadro}
Limita posloupnosti: $\forall\varepsilon\,\exists n_0:\,|a_n-L|<\varepsilon$. Racionální podíl polynomu: vyděl nejvyšší mocninou. Řada: integrální kritérium ($\sum f(n)$ konverguje $\iff \int f$ konverguje); harmonická diverguje, geometrická ($|q|<1$) $\to\frac1{1-q}$.
\end{jadro}

\kalibrace{Posloupnosti a řady jsou v požadavku analýzy a opakovaně se zkouší (geometrická řada byla samostatná otázka). Definice limity + jeden výpočet = 2 body.}
\past{Harmonická řada diverguje, ač $a_n\to0$ - samotné $a_n\to0$ je nutná, ne postačující podmínka konvergence.}
